\chapter{Cluster Analysis: Algoritmi K-Means e Clustering Gerarchico}
\label{chap:dm-kmeans-hierarchical}

In questo capitolo vengono presentati i fondamenti teorici, le formulazioni algebriche e le dinamiche operative dei due paradigmi cardine dell'apprendimento non supervisionato basato su partizionamento e gerarchizzazione:
\begin{enumerate}[noitemsep]
    \item L'\textbf{algoritmo K-Means}: formulazione analitica, dinamica di Lloyd, minimizzazione dell'errore quadratico ($\mathrm{SSE}$), dimostrazione formale di convergenza tramite discesa di coordinate, euristiche di inizializzazione (K-Means++, Bisecting K-Means) e limiti strutturali;
    \item Il \textbf{Clustering Gerarchico}: approcci agglomerativi e divisivi, interpretazione dei dendrogrammi, criteri di collegamento (\textit{linkage}: MIN, MAX, Group Average, Metodo di Ward con derivazione analitica), ricorrenza universale di Lance-Williams ed equivalenza con il Minimum Spanning Tree ($\mathrm{MST}$).
\end{enumerate}

% ====================================================================
% SEZIONE 7.1: INTRODUZIONE ALLA CLUSTER ANALYSIS PARTIZIONALE E GERARCHICA
% ====================================================================
\section{Introduzione alla Cluster Analysis Partizionale e Gerarchica}
\label{sec:dm-clustering-intro-taxonomy}

Come introdotto nel Capitolo~\ref{chap:dm-similarity-clustering-foundations} (\autoref{sec:dm-cluster-analysis-principles}), la \textbf{Cluster Analysis} persegue la partizione di un dataset $\mathcal{D} = \{\mathbf{x}_1, \dots, \mathbf{x}_n\} \subset \mathbb{R}^d$ massimizzando congiuntamente la coesione interna e la separazione esterna tra i gruppi. Sulle basi della tassonomia generale dei metodi delineata nella \autoref{sec:dm-clustering-taxonomies}, questo capitolo approfondisce i due capisaldi operativi dell'apprendimento non supervisionato: l'approccio partizionale fondato su prototipi ($K$-Means) e l'approccio gerarchico basato su matrici di prossimità e dendrogrammi.

\begin{figure}[htbp]
\centering
\begin{tikzpicture}[
    conceptBox/.style={rectangle, rounded corners=5pt, draw=navy!85!black, fill=navy!8, thick, font=\footnotesize, text width=6.1cm, inner sep=7pt},
    rootBox/.style={rectangle, rounded corners=6pt, draw=purple!85!black, fill=purple!12, very thick, font=\small, align=center, text width=11.5cm, inner sep=7pt},
    arrowStyle/.style={-Stealth, thick, draw=navy!85!black}
]
    \node[rootBox] (root) at (0, 3.2) {
        \textbf{\large PARADIGMI DI CLUSTER ANALYSIS}\\[3pt]
        \footnotesize \normalfont Metodologie di partizionamento non supervisionato dello spazio $\mathcal{D} \subset \mathbb{R}^d$
    };

    \node[conceptBox, anchor=north] (part) at (-3.4, 1.1) {
        \centering\textbf{Clustering Partizionale (K-Means)}\\[6pt]
        \raggedright
        $\bullet$ Partizione piatta in $K$ cluster disgiunti\\[3pt]
        $\bullet$ Minimizzazione inerzia globale ($\mathrm{SSE}$)\\[3pt]
        $\bullet$ Prototipi a centroide e celle di Voronoi\\[3pt]
        $\bullet$ Scalabilit\`a lineare: $\mathcal{O}(n \cdot K \cdot I \cdot d)$
    };

    \node[conceptBox, anchor=north] (hier) at (3.4, 1.1) {
        \centering\textbf{Clustering Gerarchico}\\[6pt]
        \raggedright
        $\bullet$ Tassonomia nidificata (Dendrogramma)\\[3pt]
        $\bullet$ Decisioni miopi (\textit{greedy}) irreversibili\\[3pt]
        $\bullet$ Metriche di linkage su distanze a coppie\\[3pt]
        $\bullet$ Complessit\`a sovra-quadratica: $\mathcal{O}(n^2 \log n)$
    };

    \coordinate (split) at (0, 1.9);
    \draw[thick, draw=navy!85!black] (root.south) -- (split);
    \draw[arrowStyle] (split) -| (part.north);
    \draw[arrowStyle] (split) -| (hier.north);
\end{tikzpicture}
\caption{Tassonomia dei due principali paradigmi di clustering trattati nel capitolo.}
\label{fig:dm-clustering-taxonomy-overview}
\end{figure}

I due paradigmi affrontano l'obiettivo del raggruppamento con filosofie strutturali differenti: il paradigma partizionale fissa a priori il numero di raggruppamenti $K$ e persegue una tassellazione poligonale ottima dello spazio mediante prototipi; il paradigma gerarchico genera una sequenza continua di partizioni nidificate, delegando la determinazione di $K$ a un'analisi a posteriori del dendrogramma.

% ====================================================================
% SEZIONE 7.2: K-MEANS: ALGORITMO DI LLOYD E GEOMETRIA DI VORONOI
% ====================================================================
\section{K-Means: Algoritmo di Lloyd e Geometria di Voronoi}
\label{sec:dm-kmeans-algorithm-geometry}

L'algoritmo \textbf{K-Means} rappresenta il capostipite dei metodi partizionali basati su centroide o prototipo. Introdotto storicamente da Hugo Steinhaus nel 1956, formalizzato da Stuart Lloyd (1957/1982) nei contesti di quantizzazione dei segnali, e denominato da James MacQueen (1967), il metodo realizza un partizionamento rigido (\textit{hard clustering}) in cui ogni punto è assegnato univocamente a un solo cluster.

\begin{definition}[Partizione Hard e Centroide]
Data una collezione di $n$ punti $\mathcal{D} \subset \mathbb{R}^d$ e un intero positivo fissato $K \in \mathbb{N}^+$, una partizione rigida è una famiglia di insiemi $\{C_1, \dots, C_K\}$ tale che:
\begin{equation}
    \bigcup_{i=1}^K C_i = \mathcal{D}, \qquad C_i \cap C_j = \emptyset \quad (\forall i \ne j), \qquad C_i \ne \emptyset \quad (\forall i \in \{1, \dots, K\})
\end{equation}
A ciascun cluster $C_i$ è associato un vettore prototipo denominato \textbf{centroide} $\mathbf{m}_i \in \mathbb{R}^d$ (indicato talvolta anche con $\mathbf{c}_i$), che per la distanza euclidea coincide con il baricentro aritmetico dei punti assegnati:
\begin{equation}
    \mathbf{m}_i = \frac{1}{|C_i|} \sum_{\mathbf{x} \in C_i} \mathbf{x}
    \label{eq:dm-centroid-definition}
\end{equation}
dove $|C_i|$ denota la cardinalità del cluster $C_i$.
\end{definition}

\subsection{Pseudocodice dell'Algoritmo di Lloyd}

L'algoritmo di base opera alternando ricorsivamente due fasi deterministiche fino a convergenza, come illustrato nell'Algoritmo~\ref{alg:dm-kmeans-lloyd}.

\begin{algorithm}[htbp]
\SetAlgoLined
\KwInput{Dataset $\mathcal{D} = \{\mathbf{x}_1, \dots, \mathbf{x}_n\} \subset \mathbb{R}^d$, numero di cluster $K$, tolleranza $\varepsilon \ge 0$.}
\KwOutput{Partizione finale $\{C_1, \dots, C_K\}$ e centroidi $\{\mathbf{m}_1, \dots, \mathbf{m}_K\}$.}
Inizializza $K$ centroidi iniziali $\{\mathbf{m}_1^{(0)}, \dots, \mathbf{m}_K^{(0)}\}$ mediante estrazione casuale da $\mathcal{D}$\;
$t \leftarrow 0$\;
\Repeat{i centroidi non variano significativamente: $\sum_{i=1}^K \|\mathbf{m}_i^{(t+1)} - \mathbf{m}_i^{(t)}\|^2 \le \varepsilon$}{
    $t \leftarrow t + 1$\;
    \BlankLine
    \tcp{Fase 1: Assegnazione alle celle di Voronoi}
    Inizializza $C_i^{(t)} \leftarrow \emptyset$ per ogni $i \in \{1, \dots, K\}$\;
    \ForEach{$\mathbf{x} \in \mathcal{D}$}{
        Determina l'indice del centroide più vicino: $k^* = \arg\min_{j \in \{1, \dots, K\}} \|\mathbf{x} - \mathbf{m}_j^{(t-1)}\|^2$\;
        Assegna $\mathbf{x}$ al cluster corrispondente: $C_{k^*}^{(t)} \leftarrow C_{k^*}^{(t)} \cup \{\mathbf{x}\}$\;
    }
    \BlankLine
    \tcp{Fase 2: Ricalcolo dei centroidi (baricentri)}
    \ForEach{$i \in \{1, \dots, K\}$}{
        \If{$|C_i^{(t)}| > 0$}{
            $\mathbf{m}_i^{(t)} \leftarrow \dfrac{1}{|C_i^{(t)}|} \sum_{\mathbf{x} \in C_i^{(t)}} \mathbf{x}$\;
        }
        \Else{
            Gestisci cluster vuoto riallocando un centroide orfano\;
        }
    }
}
\caption{Algoritmo standard di K-Means (Lloyd).}
\label{alg:dm-kmeans-lloyd}
\end{algorithm}

\subsection{Interpretazione Geometrica: Diagramma di Voronoi}

La fase di assegnazione partiziona lo spazio euclideo continuo $\mathbb{R}^d$ in $K$ regioni convesse delimitate da iperpiani bisettori, note come \textbf{tassellazione di Voronoi} generata dai prototipi $\{\mathbf{m}_1, \dots, \mathbf{m}_K\}$:
\begin{equation}
    V(C_i) = \left\{ \mathbf{x} \in \mathbb{R}^d : \|\mathbf{x} - \mathbf{m}_i\|_2 \le \|\mathbf{x} - \mathbf{m}_j\|_2 \quad \forall j \ne i \right\}
\end{equation}
Ciascun iperpiano di separazione tra il cluster $C_i$ e il cluster $C_j$ è ortogonale al segmento congiungente $\mathbf{m}_i$ e $\mathbf{m}_j$ e passa per il loro punto medio. Ne consegue che ogni cluster individuato da K-Means è un poliedro convesso: punti appartenenti a regioni non convesse o concave non possono essere raggruppati correttamente senza frammentare la tassellazione.

\begin{figure}[htbp]
\centering
\begin{tikzpicture}[
    stepBox/.style={rectangle, rounded corners=4pt, draw=navy!80!black, fill=navy!8, thick, font=\footnotesize, align=center, text width=6.2cm, inner sep=6pt},
    arrowStyle/.style={-Stealth, thick, draw=navy!85!black}
]
    \node[stepBox] (init) at (0, 3.6) {
        \textbf{Inizializzazione}\\
        Selezione di $K$ centroidi iniziali $\mathbf{m}_1, \dots, \mathbf{m}_K \in \mathbb{R}^d$
    };

    \node[stepBox] (assign) at (0, 1.8) {
        \textbf{Fase 1: Assegnazione (Voronoi)}\\
        $C_i = \{\mathbf{x} \in \mathcal{D} : \|\mathbf{x} - \mathbf{m}_i\|_2 \le \|\mathbf{x} - \mathbf{m}_j\|_2 \; \forall j \ne i\}$
    };

    \node[stepBox] (update) at (0, 0.0) {
        \textbf{Fase 2: Aggiornamento Prototipi}\\
        $\mathbf{m}_i = \frac{1}{|C_i|} \sum_{\mathbf{x} \in C_i} \mathbf{x} \quad (\forall i = 1, \dots, K)$
    };

    \node[draw=purple!85!black, fill=purple!12, rounded corners=4pt, font=\footnotesize\bfseries, align=center, text width=5.0cm] (stop) at (0, -1.8) {
        Convergenza raggiunta?\\
        I centroidi non variano più
    };

    \draw[arrowStyle] (init) -- (assign);
    \draw[arrowStyle] (assign) -- (update);
    \draw[arrowStyle] (update) -- (stop);
    
    \draw[arrowStyle] (stop.east) -- ++(1.8, 0) |- node[near start, right, font=\scriptsize\bfseries] {No (nuova iterazione)} (assign.east);
\end{tikzpicture}
\caption{Ciclo iterativo di aggiornamento a due fasi dell'algoritmo K-Means.}
\label{fig:dm-kmeans-flowchart}
\end{figure}

\subsection{Dinamica Temporale delle Iterazioni}

Studiando l'evoluzione dei centroidi lungo l'asse delle iterazioni $t = 1, 2, \dots$:
\begin{itemize}[noitemsep]
    \item \textbf{Iterazioni 1--3}: Si osserva la migrazione macroscopica più rilevante. I centroidi iniziali si allontanano rapidamente dalle configurazioni spurie e migrano verso i baricentri delle macro-regioni a maggiore densità.
    \item \textbf{Iterazioni 4--6 e successive}: Si entra nella fase di raffinamento fine. Pochi punti di confine vengono riassegnati lungo le superfici di Voronoi finché la variazione dei centroidi si annulla o scende sotto la soglia numerica prefissata $\varepsilon$.
\end{itemize}

\subsection{Generalizzazione ad Altre Metriche di Prossimità}

Sebbene K-Means sia formalizzato per la distanza euclidea, la struttura a due fasi può essere generalizzata ad altre misure di prossimità, a patto di ridefinire coerentemente la nozione di prototipo:
\begin{itemize}
    \item \textbf{Distanza Euclidea ($L_2$)}: Prototipo = \textbf{Media} (baricentro). Minimizza l'errore quadratico ($\mathrm{SSE}$).
    \item \textbf{Distanza di Manhattan ($L_1$)}: Prototipo = \textbf{Mediana} per componente (algoritmo \textbf{K-Medians}). Minimizza la somma degli errori assoluti ($\mathrm{SAE} = \sum_{i} \sum_{\mathbf{x} \in C_i} \|\mathbf{x} - \mathbf{m}_i\|_1$).
    \item \textbf{Similarità Coseno}: Prototipo = \textbf{Media normalizzata} a norma unitaria ($\mathbf{m}_i / \|\mathbf{m}_i\|_2$). Massimizza la similarità coseno cumulativa (impiegato nel text mining).
    \item \textbf{Correlazione di Pearson}: Prototipo = \textbf{Vettore standardizzato} (media zero, varianza unitaria).
\end{itemize}

\subsection{Analisi della Complessità Computazionale}

L'algoritmo di Lloyd vanta un'eccellente scalabilità computazionale:
\begin{itemize}
    \item \textbf{Complessità Temporale}: Ad ogni iterazione, la fase di assegnazione confronta ciascuno degli $n$ punti con tutti i $K$ centroidi in dimensione $d$, richiedendo $\mathcal{O}(n \cdot K \cdot d)$ operazioni. Il ricalcolo dei baricentri richiede la somma dei vettori, con costo $\mathcal{O}(n \cdot d)$. Indicando con $I$ il numero totale di iterazioni fino a convergenza, la complessità asintotica totale è:
    \begin{equation}
        \mathcal{T}_{\text{K-Means}} = \mathcal{O}(n \cdot K \cdot I \cdot d)
    \end{equation}
    Poiché sperimentalmente $K \ll n$, $d \ll n$ e $I \ll n$ (spesso $I \le 50$), il tempo cresce linearmente con la dimensione del dataset $n$.
    \item \textbf{Complessità Spaziale}: È sufficiente memorizzare la matrice dei dati $n \times d$ e i $K$ vettori centroide di dimensione $d$, oltre al vettore degli indici di appartenenza $n \times 1$:
    \begin{equation}
        \mathcal{S}_{\text{K-Means}} = \mathcal{O}((n + K) \cdot d)
    \end{equation}
    rendendo l'algoritmo applicabile a contesti di Big Data in cui l'intero dataset risiede in memoria principale.
\end{itemize}

% ====================================================================
% SEZIONE 7.3: FUNZIONE OBIETTIVO SSE E DIMOSTRAZIONE DI CONVERGENZA
% ====================================================================
\section{Funzione Obiettivo SSE e Dimostrazione di Convergenza}
\label{sec:dm-kmeans-sse-convergence}

L'efficacia e la stabilità analitica del K-Means derivano dall'esistenza di una funzione obiettivo scalare rigorosamente minimizzata dall'alternanza delle due fasi di Lloyd.

\begin{definition}[Sum of Squared Errors (SSE)]
Data una partizione $\{C_1, \dots, C_K\}$ e i rispettivi centroidi $\{\mathbf{m}_1, \dots, \mathbf{m}_K\}$, la \textbf{somma degli errori quadratici} ($\mathrm{SSE}$, o inerzia interna del clustering) è definita come:
\begin{equation}
    \mathrm{SSE} = \sum_{i=1}^K \sum_{\mathbf{x} \in C_i} \|\mathbf{x} - \mathbf{m}_i\|_2^2 = \sum_{i=1}^K \sum_{\mathbf{x} \in C_i} \sum_{j=1}^d (x_j - m_{ij})^2
    \label{eq:dm-sse-definition}
\end{equation}
\end{definition}

\subsection{Dimostrazione: Il Baricentro Minimizza Rigorosamente l'SSE}

\begin{theorem}[Ottimalità del Baricentro]
Dato un cluster $C_k$ contenente $|C_k|$ punti in $\mathbb{R}^d$, il prototipo $\mathbf{m}_k \in \mathbb{R}^d$ che minimizza la somma delle distanze euclidee al quadrato dei punti di $C_k$ è univocamente il baricentro aritmetico dei punti stessi.
\end{theorem}

\begin{proof}
L'errore quadratico associato al singolo cluster $C_k$ può essere scomposto come somma indipendente sulle $d$ dimensioni ortogonali:
\begin{equation}
    \mathrm{SSE}_k = \sum_{\mathbf{x} \in C_k} \sum_{j=1}^d (x_j - m_{kj})^2 = \sum_{j=1}^d \left( \sum_{\mathbf{x} \in C_k} (x_j - m_{kj})^2 \right)
\end{equation}
Per trovare il punto di minimo rispetto a ciascuna componente scalare $m_{kj}$ del centroide $\mathbf{m}_k$, calcoliamo la derivata parziale prima e la poniamo uguale a zero:
\begin{align}
    \frac{\partial \, \mathrm{SSE}_k}{\partial m_{kj}} &= \frac{\partial}{\partial m_{kj}} \sum_{\mathbf{x} \in C_k} (x_j - m_{kj})^2 \notag \\
    &= \sum_{\mathbf{x} \in C_k} 2 (x_j - m_{kj}) \cdot (-1) \notag \\
    &= -2 \sum_{\mathbf{x} \in C_k} x_j + 2 \sum_{\mathbf{x} \in C_k} m_{kj} = 0
\end{align}
Poiché $m_{kj}$ è una costante scalare rispetto all'indice di sommatoria $\mathbf{x} \in C_k$, la seconda sommatoria restituisce $|C_k| \cdot m_{kj}$:
\begin{equation}
    2 |C_k| \, m_{kj} = 2 \sum_{\mathbf{x} \in C_k} x_j \implies m_{kj} = \frac{1}{|C_k|} \sum_{\mathbf{x} \in C_k} x_j
\end{equation}
Verifichiamo la condizione di minimo calcolando la derivata seconda:
\begin{equation}
    \frac{\partial^2 \, \mathrm{SSE}_k}{\partial m_{kj}^2} = 2 |C_k| > 0
\end{equation}
La matrice Hessiana associata è una matrice diagonale scalare definita positiva:
\begin{equation}
    \nabla^2 \, \mathrm{SSE}_k = 2 |C_k| \, \mathbf{I}_d \succ 0
\end{equation}
Pertanto, il baricentro $\mathbf{m}_k = \frac{1}{|C_k|} \sum_{\mathbf{x} \in C_k} \mathbf{x}$ è l'unico minimizzatore globale stretto di $\mathrm{SSE}_k$.
\end{proof}

\subsection{Formulazione come Discesa di Coordinate e Convergenza}

L'algoritmo di Lloyd può essere formalizzato matematicamente come un'ottimizzazione per \textbf{discesa di coordinate} (\textit{Coordinate Descent}) su una funzione bivariata.

Definiamo la matrice di appartenenza binaria $\mathbf{W} \in \{0, 1\}^{n \times K}$, dove $w_{ik} = 1$ se l'istanza $\mathbf{x}_i$ appartiene al cluster $k$ e $w_{ik} = 0$ altrimenti, con il vincolo $\sum_{k=1}^K w_{ik} = 1$ per ogni $i \in \{1, \dots, n\}$. Sia inoltre $\mathbf{S} = (\mathbf{m}_1, \dots, \mathbf{m}_K) \in \mathbb{R}^{d \times K}$ l'insieme dei centroidi. La funzione obiettivo globale è:
\begin{equation}
    G(\mathbf{W}, \mathbf{S}) = \sum_{i=1}^n \sum_{k=1}^K w_{ik} \|\mathbf{x}_i - \mathbf{m}_k\|_2^2
    \label{eq:dm-kmeans-coordinate-descent-obj}
\end{equation}

\begin{theorem}[Convergenza Monotona a un Minimo Locale]
L'algoritmo K-Means riduce monotonicamente la funzione obiettivo $G(\mathbf{W}, \mathbf{S})$ ad ogni semi-passo e converge in un numero finito di iterazioni a un minimo locale o punto stazionario.
\end{theorem}

\begin{proof}
L'esecuzione si articola in due passi di minimizzazione alternata:
\begin{enumerate}
    \item \textbf{Passo 1 (Assegnazione)}: Fissata la matrice dei centroidi $\mathbf{S}^{(t-1)}$, si minimizza $G(\mathbf{W}, \mathbf{S}^{(t-1)})$ rispetto a $\mathbf{W}$. Poiché ogni termine per l'osservazione $\mathbf{x}_i$ è indipendente, la funzione è minimizzata ponendo:
    \begin{equation}
        w_{ik}^{(t)} = \begin{cases}
            1 & \text{se } k = \arg\min_{j} \|\mathbf{x}_i - \mathbf{m}_j^{(t-1)}\|_2^2 \\
            0 & \text{altrimenti}
        \end{cases}
    \end{equation}
    Ne consegue che:
    \begin{equation}
        G(\mathbf{W}^{(t)}, \mathbf{S}^{(t-1)}) \le G(\mathbf{W}^{(t-1)}, \mathbf{S}^{(t-1)})
    \end{equation}
    \item \textbf{Passo 2 (Aggiornamento)}: Fissata la nuova matrice di appartenenza $\mathbf{W}^{(t)}$, si minimizza $G(\mathbf{W}^{(t)}, \mathbf{S})$ rispetto a $\mathbf{S}$. Dal teorema precedente sull'ottimalità del baricentro, la scelta ottima per ogni $\mathbf{m}_k$ è:
    \begin{equation}
        \mathbf{m}_k^{(t)} = \frac{\sum_{i=1}^n w_{ik}^{(t)} \mathbf{x}_i}{\sum_{i=1}^n w_{ik}^{(t)}} = \frac{1}{|C_k^{(t)}|} \sum_{\mathbf{x} \in C_k^{(t)}} \mathbf{x}
    \end{equation}
    Ne consegue che:
    \begin{equation}
        G(\mathbf{W}^{(t)}, \mathbf{S}^{(t)}) \le G(\mathbf{W}^{(t)}, \mathbf{S}^{(t-1)})
    \end{equation}
\end{enumerate}
Combinando le due disuguaglianze, la successione dei valori assunti dalla funzione obiettivo è monotona non crescente:
\begin{equation}
    G(\mathbf{W}^{(t)}, \mathbf{S}^{(t)}) \le G(\mathbf{W}^{(t-1)}, \mathbf{S}^{(t-1)})
\end{equation}
Poiché $G(\mathbf{W}, \mathbf{S}) \ge 0$ è limitata inferiormente da zero, e il numero di possibili matrici di assegnazione discrete $\mathbf{W}$ è finito (pari al numero di partizioni di $n$ elementi in $K$ gruppi non vuoti, quantificato dai numeri di Stirling di seconda specie $S(n, K)$), l'algoritmo non può entrare in cicli infiniti e deve necessariamente convergere in un numero finito di passi a una configurazione stazionaria (minimo locale).
\end{proof}

% ====================================================================
% SEZIONE 7.4: SCELTA DI K, MINIMI LOCALI E INIZIALIZZAZIONE
% ====================================================================
\section{Scelta del Parametro K, Minimi Locali e Inizializzazione}
\label{sec:dm-kmeans-k-choice-local-minima}

Nonostante la convergenza garantita, la superficie di errore dell'ottimizzazione K-Means è altamente non convessa e costellata da innumerevoli minimi locali sub-ottimali. Infatti, la determinazione della partizione che minimizza globalmente l'SSE è un problema \textbf{NP-hard} anche per soli $K = 2$ cluster in dimensione arbitraria, o per dimensione $d = 2$ con $K$ arbitrario.

\subsection{Determinazione del Numero di Cluster: Il Metodo a Gomito}
\label{sec:dm-kmeans-elbow-method}

Poiché all'aumentare di $K$ la funzione $\mathrm{SSE}(K)$ decresce monotonicamente (fino ad annullarsi banalmente per $K = n$), non è possibile minimizzare l'SSE direttamente rispetto a $K$. 

Si ricorre all'\textbf{Elbow Method} (Metodo a Gomito):
\begin{enumerate}
    \item Si esegue K-Means per una sequenza crescente di valori candidati $K \in \{1, 2, \dots, K_{\max}\}$;
    \item Si traccia la curva di $\mathrm{SSE}$ in funzione di $K$;
    \item Si individua il punto di flesso ("gomito") in cui il decremento marginale di errore subisce un drastico rallentamento.
\end{enumerate}
Matematicamente, il gomito corrisponde al valore di $K$ che massimizza la derivata seconda discreta dell'SSE:
\begin{equation}
    \Delta^2 \, \mathrm{SSE}(K) = \mathrm{SSE}(K+1) - 2 \, \mathrm{SSE}(K) + \mathrm{SSE}(K-1)
\end{equation}
Oltre tale soglia, l'introduzione di ulteriori cluster frammenta aggregazioni naturali senza apportare reale guadagno informativo.

\subsection{Sensibilità ai Centroidi Iniziali e Dimostrazione Combinatoria}

La qualità della partizione finale di K-Means dipende criticamente dal posizionamento dei centroidi iniziali $\mathbf{m}_1^{(0)}, \dots, \mathbf{m}_K^{(0)}$. Se due centroidi cadono all'interno dello stesso addensamento naturale, uno dei cluster vicini rimarrà orfano, intrappolando l'algoritmo in un minimo locale sub-ottimale.

\begin{theorem}[Probabilità Combinatoria di Inizializzazione Perfetta]
Si consideri un dataset ideale composto da $K$ cluster naturali ben separati, ciascuno avente la medesima cardinalità $N/K$. Scegliendo $K$ centroidi iniziali mediante campionamento casuale uniforme indipendente con reinserimento, la probabilità $P(K)$ che ciascun centroide iniziale appartenga a un cluster distinto è:
\begin{equation}
    P(K) = \frac{K!}{K^K}
    \label{eq:dm-init-prob-formula}
\end{equation}
\end{theorem}

\begin{proof}
Calcoliamo la probabilità come prodotto delle probabilità condizionate passo per passo:
\begin{enumerate}
    \item Il primo centroide può cadere in uno qualsiasi dei $K$ cluster: $P_1 = \frac{K}{K} = 1$.
    \item Affinché il secondo centroide cada in un cluster distinto dal primo, deve essere estratto da uno dei restanti $K-1$ cluster non ancora occupati. Poiché tutti i cluster hanno cardinalità uguale, la frazione di punti ammissibili è $P_2 = \frac{K-1}{K}$.
    \item Il terzo centroide deve cadere in uno dei restanti $K-2$ cluster: $P_3 = \frac{K-2}{K}$.
    \item Iterando fino al $K$-esimo centroide, la frazione residua è $P_K = \frac{1}{K}$.
\end{enumerate}
Moltiplicando le probabilità degli eventi indipendenti:
\begin{equation}
    P(K) = \prod_{j=1}^K \frac{K - j + 1}{K} = \frac{K \cdot (K-1) \cdot (K-2) \cdots 1}{K^K} = \frac{K!}{K^K}
\end{equation}
\end{proof}

Valutando la formula~\eqref{eq:dm-init-prob-formula} per diversi valori di $K$:
\begin{itemize}[noitemsep]
    \item $K = 2 \implies P(2) = \frac{2}{4} = 0.50$ (50\%);
    \item $K = 3 \implies P(3) = \frac{6}{27} \approx 0.222$ (22.2\%);
    \item $K = 5 \implies P(5) = \frac{120}{3125} \approx 0.0384$ (3.84\%);
    \item $K = 10 \implies P(10) = \frac{3\,628\,800}{10^{10}} \approx 0.00036$ (appena lo \textbf{0.036\%}).
\end{itemize}

Applicando l'approssimazione asintotica di Stirling $K! \approx \sqrt{2\pi K} \left(\frac{K}{e}\right)^K$, otteniamo:
\begin{equation}
    P(K) \approx \frac{\sqrt{2\pi K} \left(\frac{K}{e}\right)^K}{K^K} = \sqrt{2\pi K} \cdot e^{-K}
\end{equation}
La probabilità di successo dell'inizializzazione casuale pura decade in modo \textbf{esponenzialmente veloce} verso zero al crescere di $K$. Con $K \ge 10$, l'inizializzazione uniforme non garantisce quasi mai la corretta convergenza.

% ====================================================================
% SEZIONE 7.5: EURISTICHE DI INIZIALIZZAZIONE E VARIANTI DI K-MEANS
% ====================================================================
\section{Euristiche di Inizializzazione e Varianti di K-Means}
\label{sec:dm-kmeans-heuristics-variants}

Per superare la fragilità dell'inizializzazione casuale, la letteratura ha sviluppato potenti strategie correttive ed euristiche algoritmiche avanzate.

\subsection{Tecniche Euristiche di Inizializzazione}

\begin{enumerate}
    \item \textbf{Restart Casuali Multipli (\textit{Random Restarts})}: Si esegue l'algoritmo $R$ volte (es. $R = 50$) con semi casuali differenti, selezionando la soluzione che restituisce il valore minimo di $\mathrm{SSE}$. Efficace per valori modesti di $K$, ma computazionalmente oneroso e inefficace per $K$ elevato.
    \item \textbf{Campionamento con Pre-clustering Gerarchico}: Si estrae un campione casuale ridotto del dataset ($\approx 1\% - 5\%$) e si applica un algoritmo gerarchico agglomerativo con Metodo di Ward per determinare $K$ raggruppamenti stabili, i cui baricentri vengono usati come semi per K-Means.
    \item \textbf{K-Means++ ($D^2$-Weighting)}: Introdotto da Arthur e Vassilvitskii (2007), sceglie il primo centroide a caso, e ciascun centroide successivo con probabilità proporzionale al quadrato della distanza minima dai centroidi già allocati:
    \begin{equation}
        p(\mathbf{x}) = \frac{D(\mathbf{x})^2}{\sum_{\mathbf{y} \in \mathcal{D}} D(\mathbf{y})^2}, \qquad D(\mathbf{x}) = \min_{j \in \{1, \dots, m\}} \|\mathbf{x} - \mathbf{m}_j\|_2
    \end{equation}
    \textit{Garanzia teorica}: K-Means++ garantisce che il valore atteso dell'SSE finale sia entro un fattore logaritmico rispetto all'ottimo globale:
    \begin{equation}
        \mathbb{E}[\mathrm{SSE}] \le 8(\ln K + 2) \cdot \mathrm{SSE}_{\text{opt}} = \mathcal{O}(\log K) \cdot \mathrm{SSE}_{\text{opt}}
    \end{equation}
    \item \textbf{Bisecting K-Means}: Algoritmo partizionale ibrido che opera ricorsivamente con bisezioni binarie ($K=2$):
    \begin{itemize}[noitemsep]
        \item Inizia considerando l'intero dataset come un unico macro-cluster;
        \item A ogni iterazione, seleziona il cluster con il più alto $\mathrm{SSE}$ interno (o con cardinalità massima);
        \item Applica il K-Means standard con $K=2$ su tale cluster per scinderlo in due sotto-gruppi;
        \item Ripete il processo finché non si ottengono esattamente $K$ cluster.
    \end{itemize}
    Bisecting K-Means è nettamente meno sensibile all'inizializzazione rispetto a Lloyd e produce naturalmente una gerarchia parziale dei cluster.
\end{enumerate}

\subsection{Aggiornamento Incrementale (Online K-Means)}

Nella variante di MacQueen (1967), i centroidi vengono aggiornati \textit{immediatamente} dopo la riassegnazione di ogni singolo punto, invece di attendere la fine dell'epoca (\textit{batch}):
\begin{property}[Formule di Aggiornamento Incrementale di MacQueen]
Se il punto $\mathbf{x}$ viene rimosso dal cluster $C_j$ e assegnato al cluster $C_i$:
\begin{align}
    \mathbf{m}_j &\leftarrow \mathbf{m}_j - \frac{\mathbf{x} - \mathbf{m}_j}{|C_j| - 1} \label{eq:dm-macqueen-remove} \\
    \mathbf{m}_i &\leftarrow \mathbf{m}_i + \frac{\mathbf{x} - \mathbf{m}_i}{|C_i| + 1} \label{eq:dm-macqueen-add}
\end{align}
\end{property}
L'approccio incrementale garantisce che l'SSE decresca strettamente ad ogni singolo spostamento, azzera la possibilità che si generino cluster vuoti, ma introduce una dipendenza dall'ordine di presentazione dei punti.

\subsection{Gestione dei Cluster Vuoti (\textit{Empty Clusters})}

Nell'algoritmo batch di Lloyd, può accadere che durante la fase di assegnazione nessun punto dati sia più vicino al centroide $\mathbf{m}_k$ rispetto agli altri $K-1$ centroidi, portando a $|C_k| = 0$ (centroide orfano).
Strategie di ripristino per mantenere inalterato $K$:
\begin{enumerate}
    \item \textbf{Selezione a Massimo Errore Residuo}: Identificare il punto $\mathbf{x}^* \in \mathcal{D}$ che massimizza la distanza quadratica dal proprio centroide assegnato $\|\mathbf{x}^* - \mathbf{m}_{c(\mathbf{x}^*)}\|^2$ e promuoverlo a nuovo centroide isolato del cluster vuoto. Questa strategia produce la massima riduzione marginale istantanea dell'SSE.
    \item \textbf{Scissione del Cluster a Massimo SSE}: Identificare il cluster che esibisce il più elevato $\mathrm{SSE}$ cumulativo ed estrarre il suo punto più periferico per popolare il cluster orfano.
\end{enumerate}

% ====================================================================
% SEZIONE 7.6: LIMITAZIONI STRUTTURALI DI K-MEANS E TECNICHE DI ADATTAMENTO
% ====================================================================
\section{Limitazioni Strutturali di K-Means e Tecniche di Adattamento}
\label{sec:dm-kmeans-limitations-isodata}

K-Means presuppone implicitamente che i cluster siano sferici, compatti, ben separati e dotati di volumi e densità omogenee. In presenza di configurazioni geometriche complesse, l'algoritmo fallisce sistematicamente.

\begin{figure}[htbp]
\centering
\begin{tikzpicture}[
    failCard/.style={rectangle, rounded corners=5pt, draw=crimson!85!black, fill=crimson!6, thick, font=\footnotesize, align=left, text width=6.0cm, inner sep=6pt}
]
    \node[failCard] (c1) at (-3.4, 2.0) {
        \textbf{1. Dimensioni Eterogenee}\\
        I centroidi dividono equamente lo spazio; il cluster più piccolo attrae e ingloba punti periferici del cluster più grande.
    };

    \node[failCard] (c2) at (3.4, 2.0) {
        \textbf{2. Densità Eterogenee}\\
        I cluster sparsi a bassa densità vengono spezzati, mentre cluster compatti ad alta densità vicini vengono fusi.
    };

    \node[failCard] (c3) at (-3.4, -0.4) {
        \textbf{3. Forme Non Globulari}\\
        Geometrie concave o intrecciate (\textit{two moons}, anelli) vengono tagliate linearmente dagli iperpiani di Voronoi.
    };

    \node[failCard] (c4) at (3.4, -0.4) {
        \textbf{4. Sensibilità agli Outlier}\\
        La metrica quadratica $\|\mathbf{x} - \mathbf{m}\|^2$ devia gravemente i centroidi verso le anomalie isolate.
    };
\end{tikzpicture}
\caption{Quattro scenari critici di fallimento strutturale dell'algoritmo standard K-Means.}
\label{fig:dm-kmeans-structural-limitations}
\end{figure}

\subsection{Analisi delle Tre Limitazioni Geometriche}

\begin{enumerate}
    \item \textbf{Dimensioni Differenti}: Poiché l'SSE penalizza gli scarti quadratici senza considerare il volume intrinseco, il centroide assegnato al cluster compatto attrae punti del cluster disperso per minimizzare l'errore globale, determinando confini incoerenti.
    \item \textbf{Densità Differenti}: Iperpiani posizionati equidistantemente tra centroidi non tengono conto del gradiente di densità. Un cluster a bassissima densità richiederà più centroidi per essere coperto, causando l'unione impropria di cluster ad alta densità vicini.
    \item \textbf{Forme Non Globulari (es. \textit{Two Moons})}: Le celle di Voronoi sono poliedri convessi. Se i dati seguono varietà curvilinee non convesse (due semicerchi intrecciati), un iperpiano lineare secante taglia inevitabilmente a metà entrambi i corpi continui, fallendo completamente la separazione.
\end{enumerate}

\begin{noteblock}{Strategia di Risoluzione: Over-Clustering}
Per modellare varietà curvilinee o densità eterogenee, si esegue K-Means con un numero elevato di micro-cluster ($K' \gg K$). Tanti piccoli centroidi formano un'approssimazione poligonale a pezzi (\textit{piecewise linear approximation}) che segue il profilo della curva continua. Successivamente, algoritmi basati su grafi o collegamenti ricompongono i micro-cluster nelle forme non lineari desiderate.
\end{noteblock}

\subsection{Pre-processing e Post-processing: Il Paradigma ISODATA}

Per mitigare questi limiti nell'analisi partizionale, si adottano routine di pre- e post-elaborazione formalizzate storicamente nell'algoritmo \textbf{ISODATA} (\textit{Iterative Self-Organizing Data Analysis Technique}):
\begin{itemize}
    \item \textbf{Pre-processing}:
    \begin{itemize}[noitemsep]
        \item Normalizzazione o standardizzazione delle feature (z-score) per impedire che dimensioni con scale numeriche elevate dominino il calcolo della distanza euclidea;
        \item Identificazione e rimozione preliminare degli outlier prima dell'inizializzazione dei centroidi.
    \end{itemize}
    \item \textbf{Post-processing (Dinamica ISODATA)}:
    \begin{itemize}[noitemsep]
        \item \textbf{Split (Scissione)}: Se la deviazione standard o l'SSE di un cluster supera una soglia massima $\theta_{\text{split}}$, e la cardinalità è sufficiente ($|C_i| \ge 2\theta_{\min}$), il cluster viene diviso in due lungo la coordinata di massima varianza.
        \item \textbf{Merge (Fusione)}: Se la distanza euclidea tra due centroidi $\|\mathbf{m}_i - \mathbf{m}_j\|_2$ è inferiore a una soglia minima $\theta_{\text{merge}}$, i due cluster vengono fusi in un unico raggruppamento baricentrico.
        \item \textbf{Eliminazione}: Se $|C_i| < \theta_{\min}$, il cluster viene soppresso e i suoi punti redistribuiti.
    \end{itemize}
\end{itemize}

% ====================================================================
% SEZIONE 7.7: FONDAMENTI DEL CLUSTERING GERARCHICO
% ====================================================================
\section{Fondamenti del Clustering Gerarchico}
\label{sec:dm-hierarchical-foundations-dendrogram}

Il \textbf{Clustering Gerarchico} supera la rigidità della partizione piatta generando una tassonomia continua di partizioni nidificate (\textit{nested clusters}).

\begin{definition}[Dendrogramma]
Il \textbf{dendrogramma} è una rappresentazione grafica ad albero binario radicato in cui:
\begin{itemize}[noitemsep]
    \item Le \textbf{foglie} al livello $h = 0$ rappresentano i singoli punti dati isolati (cluster singleton);
    \item I \textbf{nodi interni} rappresentano le unioni successive di coppie di cluster;
    \item L'\textbf{altezza} $h$ a cui due rami si collegano corrisponde formalmente alla distanza o dissimilarità inter-cluster al momento della loro fusione.
\end{itemize}
\end{definition}

\begin{figure}[htbp]
\centering
\begin{tikzpicture}[
    every node/.style={font=\footnotesize},
    pointNode/.style={circle, fill=navy!85!black, inner sep=1.5pt},
    cutLine/.style={dashed, thick, draw=crimson!85!black}
]
    % Assi del Dendrogramma
    \draw[-Stealth, thick] (0,0) -- (0, 3.8) node[left] {Altezza $h$};
    \draw[thick] (0,0) -- (7.5, 0);

    % Foglie
    \node[below] at (0.8, 0) {$p_3$};
    \node[below] at (1.8, 0) {$p_6$};
    \node[below] at (2.8, 0) {$p_4$};
    \node[below] at (4.2, 0) {$p_2$};
    \node[below] at (5.2, 0) {$p_5$};
    \node[below] at (6.6, 0) {$p_1$};

    % 1. Fusione {p3, p6} a h = 0.8
    \draw[thick] (0.8, 0) -- (0.8, 0.8) -- (1.8, 0.8) -- (1.8, 0);
    \coordinate (c36) at (1.3, 0.8);

    % 2. Fusione {p2, p5} a h = 1.1
    \draw[thick] (4.2, 0) -- (4.2, 1.1) -- (5.2, 1.1) -- (5.2, 0);
    \coordinate (c25) at (4.7, 1.1);

    % 3. Fusione {{p3, p6}, p4} a h = 1.5
    \draw[thick] (c36) -- (1.3, 1.5) -- (2.8, 1.5) -- (2.8, 0);
    \coordinate (c364) at (2.05, 1.5);

    % 4. Fusione {{p3, p6, p4}, {p2, p5}} a h = 2.2
    \draw[thick] (c364) -- (2.05, 2.2) -- (4.7, 2.2) -- (c25);
    \coordinate (c_left) at (3.375, 2.2);

    % 5. Fusione finale con p1 a h = 3.2
    \draw[thick] (c_left) -- (3.375, 3.2) -- (6.6, 3.2) -- (6.6, 0);

    % Taglio orizzontale a quota h* = 1.8
    \draw[cutLine] (-0.2, 1.8) -- (7.5, 1.8);
    \node[left, crimson!85!black, font=\scriptsize] at (-0.05, 1.8) {$h^* = 1.8$};
    \node[right, crimson!85!black, font=\scriptsize\bfseries] at (7.5, 1.8) {Taglio $h^* \implies K = 3$};
\end{tikzpicture}
\caption{Struttura di un dendrogramma: le altezze delle giunzioni indicano la distanza di unione. Un taglio orizzontale induce una partizione piatta con $K = 3$ cluster.}
\label{fig:dm-dendrogram-schema}
\end{figure}

\subsection{Vantaggi Distintivi del Paradigma Gerarchico}

\begin{itemize}
    \item \textbf{Nessun vincolo a priori su $K$}: Non è richiesto specificare anticipatamente il numero di cluster. L'analista ottiene una visualizzazione completa dell'architettura dei dati ed estrae $K$ a posteriori tracciando un taglio orizzontale alla quota desiderata $h^*$.
    \item \textbf{Modellazione di gerarchie biologiche e tassonomiche}: Strutture evolutive, tassonomie filogenetiche e relazioni semantiche possiedono una natura intrinsecamente gerarchica perfettamente rispecchiata dal dendrogramma.
    \item \textbf{Determinismo}: A differenza del K-Means, gli algoritmi gerarchici standard sono completamente deterministici (non dipendono da semi di inizializzazione casuali).
\end{itemize}

\subsection{Tassonomia degli Algoritmi Gerarchici}

\begin{itemize}
    \item \textbf{Agglomerativi (\textit{Bottom-up})}: Iniziano considerando ogni punto come un cluster singleton ($n$ cluster iniziali) e, a ciascun passo, fondono iterativamente la coppia di cluster più vicina fino a formare un unico grande raggruppamento radice in $n - 1$ passi.
    \item \textbf{Divisivi (\textit{Top-down})}: Iniziano considerando l'intero dataset come un unico mega-cluster e, ad ogni iterazione, partizionano ricorsivamente i gruppi in sotto-gruppi fino a ricondursi a $n$ singleton.
\end{itemize}

% ====================================================================
% SEZIONE 7.8: ALGORITMO AGGLOMERATIVO E CRITERI DI COLLEGAMENTO (LINKAGE)
% ====================================================================
\section{Algoritmo Agglomerativo e Criteri di Collegamento (Linkage)}
\label{sec:dm-agglomerative-linkage-criteria}

L'approccio agglomerativo standard richiede la definizione di una matrice di prossimità inter-cluster $\mathbf{D} \in \mathbb{R}^{N \times N}$ e segue il ciclo descritto nell'Algoritmo~\ref{alg:dm-hierarchical-agglomerative}.

\begin{algorithm}[htbp]
\SetAlgoLined
\KwInput{Dataset $\mathcal{D} = \{\mathbf{x}_1, \dots, \mathbf{x}_n\}$, metrica di distanza tra punti $d(\cdot, \cdot)$.}
\KwOutput{Dendrogramma completo delle fusioni gerarchiche.}
Inizializza $n$ cluster singleton: $C_i = \{\mathbf{x}_i\}$ per ogni $i \in \{1, \dots, n\}$\;
Costruisci la matrice di prossimità iniziale $\mathbf{D}^{(0)} \in \mathbb{R}^{n \times n}$ con $D_{ij} = d(\mathbf{x}_i, \mathbf{x}_j)$\;
\For{$t = 1$ \KwTo $n - 1$}{
    Trova la coppia di cluster $(A, B)$ a minima distanza nella matrice: $(A, B) = \arg\min_{i \ne j} D_{ij}^{(t-1)}$\;
    Fondi i cluster $A$ e $B$ creando il nuovo cluster $A \cup B$ a quota $h_t = D_{AB}^{(t-1)}$\;
    Elimina le righe e colonne corrispondenti ad $A$ e $B$ dalla matrice di prossimità\;
    Calcola la distanza tra il nuovo cluster $A \cup B$ e ogni altro cluster residuo $Q$ secondo il criterio di collegamento prescelto\;
    Aggiorna la matrice di prossimità $\mathbf{D}^{(t)}$ inserendo la nuova riga/colonna associata ad $A \cup B$\;
}
\caption{Algoritmo gerarchico agglomerativo standard.}
\label{alg:dm-hierarchical-agglomerative}
\end{algorithm}

\subsection{I Quattro Criteri di Collegamento (\textit{Linkage})}

La definizione della distanza $D(A, B)$ tra due insiemi di punti $A$ e $B$ determina interamente la forma e le proprietà della gerarchia risultante.

\subsubsection{1. MIN (Single Link)}
La distanza tra due cluster coincide con la minima distanza tra qualsiasi coppia di punti:
\begin{equation}
    D_{\text{MIN}}(A, B) = \min_{\mathbf{x} \in A, \, \mathbf{y} \in B} d(\mathbf{x}, \mathbf{y})
    \label{eq:dm-single-link}
\end{equation}
\begin{itemize}[noitemsep]
    \item \textbf{Punti di forza}: Capace di catturare forme arbitrarie non convesse, filiformi o ad anello (è isomorfo alla costruzione del Minimum Spanning Tree).
    \item \textbf{Debolezze (\textit{Chaining Effect})}: Estremamente sensibile al rumore e agli outlier. Anche un singolo punto spurio situato tra due cluster compatti distinti funge da "ponte", inducendo la fusione prematura dei due macro-gruppi in un'unica lunga catena.
\end{itemize}

\subsubsection{2. MAX (Complete Link)}
La distanza tra due cluster coincide con la massima distanza tra le loro coppie di punti:
\begin{equation}
    D_{\text{MAX}}(A, B) = \max_{\mathbf{x} \in A, \, \mathbf{y} \in B} d(\mathbf{x}, \mathbf{y})
    \label{eq:dm-complete-link}
\end{equation}
\begin{itemize}[noitemsep]
    \item \textbf{Punti di forza}: Robusto rispetto al rumore e immune al chaining effect; tende a produrre cluster iper-compatti e globulari limitando formalmente il diametro massimo del cluster combinato: $\text{diam}(A \cup B) \le h$.
    \item \textbf{Debolezze}: Tende a frammentare ingiustificatamente grandi cluster allungati e soffre la presenza di outlier periferici estremi.
\end{itemize}

\subsubsection{3. Group Average (Average Link / UPGMA)}
La distanza tra due cluster è la media aritmetica di tutte le distanze a coppie:
\begin{equation}
    D_{\text{AVG}}(A, B) = \frac{1}{|A| \cdot |B|} \sum_{\mathbf{x} \in A} \sum_{\mathbf{y} \in B} d(\mathbf{x}, \mathbf{y})
    \label{eq:dm-average-link}
\end{equation}
Costituisce un compromesso matematicamente bilanciato tra la permissività di MIN e la severità di MAX. È relativamente robusto al rumore poiché l'effetto di singoli outlier viene mediato sull'intera cardinalità del gruppo.

\subsubsection{4. Metodo di Ward (Minimizzazione della Varianza)}
Introdotto da Joe H. Ward Jr. (1963), la prossimità tra due cluster è misurata tramite l'\textbf{incremento dell'errore quadratico totale} ($\Delta \mathrm{SSE}$) che deriverebbe dalla loro eventuale fusione:
\begin{equation}
    \Delta \mathrm{SSE}_{AB} = \mathrm{SSE}_{A \cup B} - (\mathrm{SSE}_A + \mathrm{SSE}_B)
    \label{eq:dm-ward-delta-def}
\end{equation}

\begin{theorem}[Formulazione Chiusa del Metodo di Ward]
L'incremento dell'errore quadratico $\Delta \mathrm{SSE}_{AB}$ dovuto alla fusione dei cluster $A$ e $B$ equivale alla distanza euclidea al quadrato tra i loro rispettivi centroidi, scalata dal prodotto armonico delle loro masse:
\begin{equation}
    \Delta \mathrm{SSE}_{AB} = \frac{|A| \cdot |B|}{|A| + |B|} \|\mathbf{m}_A - \mathbf{m}_B\|_2^2
    \label{eq:dm-ward-formula-closed}
\end{equation}
\end{theorem}

\begin{proof}
Siano $A$ e $B$ due cluster disgiunti con cardinalità $|A|, |B|$ e centroidi $\mathbf{m}_A, \mathbf{m}_B$. Il centroide del cluster fuso $A \cup B$ è la media pesata dei baricentri:
\begin{equation}
    \mathbf{m}_{A \cup B} = \frac{|A| \mathbf{m}_A + |B| \mathbf{m}_B}{|A| + |B|}
\end{equation}
Calcoliamo lo scarto vettoriale tra il baricentro combinato e il baricentro originale $\mathbf{m}_A$:
\begin{equation}
    \mathbf{m}_{A \cup B} - \mathbf{m}_A = \frac{|A|\mathbf{m}_A + |B|\mathbf{m}_B - (|A|+|B|)\mathbf{m}_A}{|A| + |B|} = \frac{|B|(\mathbf{m}_B - \mathbf{m}_A)}{|A| + |B|}
\end{equation}
Simmetricamente per il cluster $B$:
\begin{equation}
    \mathbf{m}_{A \cup B} - \mathbf{m}_B = \frac{|A|(\mathbf{m}_A - \mathbf{m}_B)}{|A| + |B|}
\end{equation}
L'errore quadratico totale del cluster unificato è:
\begin{equation}
    \mathrm{SSE}_{A \cup B} = \sum_{\mathbf{x} \in A} \|\mathbf{x} - \mathbf{m}_{A \cup B}\|^2 + \sum_{\mathbf{y} \in B} \|\mathbf{y} - \mathbf{m}_{A \cup B}\|^2
\end{equation}
Sviluppiamo il primo termine inserendo e sottraendo il baricentro originale $\mathbf{m}_A$:
\begin{align}
    \sum_{\mathbf{x} \in A} \|\mathbf{x} - \mathbf{m}_{A \cup B}\|^2 &= \sum_{\mathbf{x} \in A} \|(\mathbf{x} - \mathbf{m}_A) - (\mathbf{m}_{A \cup B} - \mathbf{m}_A)\|^2 \notag \\
    &= \sum_{\mathbf{x} \in A} \|\mathbf{x} - \mathbf{m}_A\|^2 + \sum_{\mathbf{x} \in A} \|\mathbf{m}_{A \cup B} - \mathbf{m}_A\|^2 \notag \\
    &\quad - 2 \left( \sum_{\mathbf{x} \in A} (\mathbf{x} - \mathbf{m}_A) \right) \cdot (\mathbf{m}_{A \cup B} - \mathbf{m}_A)
\end{align}
Poiché per definizione di centroide $\sum_{\mathbf{x} \in A} (\mathbf{x} - \mathbf{m}_A) = \mathbf{0}$, il doppio prodotto svanisce identicamente:
\begin{equation}
    \sum_{\mathbf{x} \in A} \|\mathbf{x} - \mathbf{m}_{A \cup B}\|^2 = \mathrm{SSE}_A + |A| \cdot \|\mathbf{m}_{A \cup B} - \mathbf{m}_A\|^2
\end{equation}
Sostituendo l'espressione precedentemente ricavata per $(\mathbf{m}_{A \cup B} - \mathbf{m}_A)$:
\begin{align}
    \sum_{\mathbf{x} \in A} \|\mathbf{x} - \mathbf{m}_{A \cup B}\|^2 &= \mathrm{SSE}_A + |A| \left\| \frac{|B|(\mathbf{m}_B - \mathbf{m}_A)}{|A| + |B|} \right\|^2 \notag \\
    &= \mathrm{SSE}_A + \frac{|A| \cdot |B|^2}{(|A| + |B|)^2} \|\mathbf{m}_A - \mathbf{m}_B\|^2
\end{align}
Ripetendo lo stesso sviluppo per il secondo termine relativo a $B$:
\begin{equation}
    \sum_{\mathbf{y} \in B} \|\mathbf{y} - \mathbf{m}_{A \cup B}\|^2 = \mathrm{SSE}_B + \frac{|A|^2 \cdot |B|}{(|A| + |B|)^2} \|\mathbf{m}_A - \mathbf{m}_B\|^2
\end{equation}
Sommando i due termini per ottenere $\mathrm{SSE}_{A \cup B}$:
\begin{align}
    \mathrm{SSE}_{A \cup B} &= \mathrm{SSE}_A + \mathrm{SSE}_B + \left[ \frac{|A| \cdot |B|^2 + |A|^2 \cdot |B|}{(|A| + |B|)^2} \right] \|\mathbf{m}_A - \mathbf{m}_B\|^2 \notag \\
    &= \mathrm{SSE}_A + \mathrm{SSE}_B + \left[ \frac{|A| \cdot |B| (|B| + |A|)}{(|A| + |B|)^2} \right] \|\mathbf{m}_A - \mathbf{m}_B\|^2 \notag \\
    &= \mathrm{SSE}_A + \mathrm{SSE}_B + \frac{|A| \cdot |B|}{(|A| + |B|)} \|\mathbf{m}_A - \mathbf{m}_B\|^2
\end{align}
Sottraendo $(\mathrm{SSE}_A + \mathrm{SSE}_B)$, si ottiene la tesi:
\begin{equation}
    \Delta \mathrm{SSE}_{AB} = \mathrm{SSE}_{A \cup B} - (\mathrm{SSE}_A + \mathrm{SSE}_B) = \frac{|A| \cdot |B|}{|A| + |B|} \|\mathbf{m}_A - \mathbf{m}_B\|^2
\end{equation}
\end{proof}

Il Metodo di Ward è l'analogo gerarchico esatto del criterio di costo del K-Means: genera cluster regolari, sferici e ben bilanciati, risultando straordinariamente robusto rispetto al rumore.

\subsection{Esempio Pratico Guidato su 6 Punti: MIN vs MAX}

Consideriamo un dataset di 6 punti bidimensionali $\{p_1, \dots, p_6\}$ con matrice di distanza iniziale:
\begin{equation}
\mathbf{D} = 
\begin{pmatrix}
 & p_1 & p_2 & p_3 & p_4 & p_5 & p_6 \\
p_1 & 0.00 & 0.24 & 0.22 & 0.37 & 0.34 & 0.23 \\
p_2 & 0.24 & 0.00 & 0.15 & 0.20 & 0.14 & 0.25 \\
p_3 & 0.22 & 0.15 & 0.00 & 0.15 & 0.28 & 0.11 \\
p_4 & 0.37 & 0.20 & 0.15 & 0.00 & 0.29 & 0.22 \\
p_5 & 0.34 & 0.14 & 0.28 & 0.29 & 0.00 & 0.39 \\
p_6 & 0.23 & 0.25 & 0.11 & 0.22 & 0.39 & 0.00
\end{pmatrix}
\end{equation}

\subsubsection{Risoluzione con MIN (Single Link)}
\begin{enumerate}
    \item Minimo assoluto globale: $d(p_3, p_6) = 0.11 \implies$ Unione di $\{p_3, p_6\}$ a quota $h = 0.11$.
    \item Minimo successivo: $d(p_2, p_5) = 0.14 \implies$ Unione di $\{p_2, p_5\}$ a quota $h = 0.14$.
    \item Distanza tra $\{p_3, p_6\}$ e $p_4$: $\min(d(p_3, p_4), d(p_6, p_4)) = 0.15 \implies$ Fusione di $\{p_3, p_6\}$ con $p_4$ a quota $h = 0.15$.
    \item Distanza tra $\{p_3, p_6, p_4\}$ e $\{p_2, p_5\}$:
    \begin{align}
        \min_{\mathbf{x} \in \{3,6,4\}, \, \mathbf{y} \in \{2,5\}} d(\mathbf{x}, \mathbf{y}) &= \min(0.15, 0.25, 0.20, 0.28, 0.39, 0.29) \notag \\
        &= d(p_3, p_2) = 0.15
    \end{align}
    I due blocchi si fondono istantaneamente a quota $h = 0.15$ sfruttando il ponte ravvicinato $p_3$--$p_2$.
    \item Distanza da $p_1$: $\min(d(p_1, p_j)) = d(p_1, p_3) = 0.22 \implies$ Fusione finale a quota $h = 0.22$.
\end{enumerate}

\subsubsection{Risoluzione con MAX (Complete Link)}
\begin{enumerate}
    \item Primi due passi identici: fusione $\{p_3, p_6\}$ a $h = 0.11$, e $\{p_2, p_5\}$ a $h = 0.14$.
    \item Distanza MAX tra $\{p_3, p_6\}$ e $p_4$: $\max(d(p_3, p_4), d(p_6, p_4)) = 0.22$.
    \item Distanza MAX tra $\{p_3, p_6\}$ e $\{p_2, p_5\}$:
    \begin{align}
        \max_{\mathbf{x} \in \{3,6\}, \, \mathbf{y} \in \{2,5\}} d(\mathbf{x}, \mathbf{y}) &= \max(0.15, 0.25, 0.28, 0.39) \notag \\
        &= d(p_6, p_5) = 0.39
    \end{align}
    Poiché $0.22 < 0.39$, l'unione preferita a questo passo è $\{p_3, p_6\}$ con $p_4$ a quota $h = 0.22$.
    \item La distanza tra $\{p_3, p_6, p_4\}$ e $\{p_2, p_5\}$ rimane elevata ($0.39$). Invece, la distanza da $p_1$ a $\{p_2, p_5\}$ è $\max(0.24, 0.34) = 0.34$. Pertanto, MAX ritarda la fusione tra i due gruppi principali ed evita fusioni precoci indotte da collegamenti isolati.
\end{enumerate}

\subsection{Formula di Ricorrenza di Lance-Williams}

La formula di ricorrenza proposta da Lance e Williams (1967) fornisce un quadro analitico unificato per aggiornare la matrice di prossimità ad ogni fusione senza dover riesaminare i singoli punti dati originari.

\begin{theorem}[Formula di Lance-Williams]
Dati due cluster $A$ e $B$ fusi nel nuovo cluster $A \cup B$, la distanza tra $A \cup B$ e un terzo cluster $Q$ è determinata da:
\begin{equation}
    D(A \cup B, \, Q) = \alpha_A \, D(A, Q) + \alpha_B \, D(B, Q) + \beta \, D(A, B) + \gamma \, |D(A, Q) - D(B, Q)|
    \label{eq:dm-lance-williams-general}
\end{equation}
dove i parametri $(\alpha_A, \alpha_B, \beta, \gamma)$ dipendono univocamente dal criterio prescelto.
\end{theorem}

\begin{table}[htbp]
\centering
\footnotesize
\setlength{\tabcolsep}{4pt}
\renewcommand{\arraystretch}{1.25}
\begin{tabular}{l c c c c}
\toprule
\textbf{Criterio di Linkage} & $\mathbf{\alpha_A}$ & $\mathbf{\alpha_B}$ & $\mathbf{\beta}$ & $\mathbf{\gamma}$ \\
\midrule
\textbf{MIN (Single Link)} & $1/2$ & $1/2$ & $0$ & $-1/2$ \\
\textbf{MAX (Complete Link)} & $1/2$ & $1/2$ & $0$ & $+1/2$ \\
\textbf{Group Average} & $\frac{|A|}{|A| + |B|}$ & $\frac{|B|}{|A| + |B|}$ & $0$ & $0$ \\[3pt]
\textbf{Centroid Linkage} & $\frac{|A|}{|A| + |B|}$ & $\frac{|B|}{|A| + |B|}$ & $-\frac{|A||B|}{(|A| + |B|)^2}$ & $0$ \\[3pt]
\textbf{Metodo di Ward} & $\frac{|A| + |Q|}{|A| + |B| + |Q|}$ & $\frac{|B| + |Q|}{|A| + |B| + |Q|}$ & $-\frac{|Q|}{|A| + |B| + |Q|}$ & $0$ \\
\bottomrule
\end{tabular}
\caption{Parametri della formula di ricorrenza di Lance-Williams per i principali metodi gerarchici.}
\label{tab:dm-lance-williams-parameters}
\end{table}

\begin{proof}[Dimostrazione di Consistenza per Group Average]
Verifichiamo che la sostituzione dei parametri nella formula di Lance-Williams riproduca esattamente la definizione originaria della media:
\begin{align}
    \sum_{\mathbf{x} \in A \cup B} \sum_{\mathbf{z} \in Q} d(\mathbf{x}, \mathbf{z}) &= \sum_{\mathbf{x} \in A} \sum_{\mathbf{z} \in Q} d(\mathbf{x}, \mathbf{z}) + \sum_{\mathbf{y} \in B} \sum_{\mathbf{z} \in Q} d(\mathbf{y}, \mathbf{z}) \notag \\
    &= |A| \cdot |Q| \cdot D_{\text{AVG}}(A, Q) + |B| \cdot |Q| \cdot D_{\text{AVG}}(B, Q)
\end{align}
Dividendo per $|A \cup B| \cdot |Q| = (|A| + |B|) \cdot |Q|$ e semplificando $|Q| > 0$:
\begin{equation}
    D_{\text{AVG}}(A \cup B, \, Q) = \frac{|A|}{|A| + |B|} D_{\text{AVG}}(A, Q) + \frac{|B|}{|A| + |B|} D_{\text{AVG}}(B, Q)
\end{equation}
che coincide esattamente con la riga di Group Average in Tabella~\ref{tab:dm-lance-williams-parameters}.
\end{proof}

\subsection{Tabella Sinottica Comparativa dei Criteri di Collegamento}

\begin{table}[htbp]
\centering
\footnotesize
\renewcommand{\arraystretch}{1.25}
\begin{tabularx}{\textwidth}{p{2.3cm} p{2.2cm} Y Y Y}
\toprule
\textbf{Criterio} & \textbf{Obiettivo} & \textbf{Geometria} & \textbf{Outlier} & \textbf{Proprietà} \\
\midrule
\textbf{MIN (Single)} & $\min_{\mathbf{x}, \mathbf{y}} d(\mathbf{x}, \mathbf{y})$ & Arbitraria, concatenata & Altissima (\textit{chaining}) & Isomorfo all'$\mathrm{MST}$ \\
\textbf{MAX (Complete)} & $\max_{\mathbf{x}, \mathbf{y}} d(\mathbf{x}, \mathbf{y})$ & Sferica, compatta & Bassa al centro & Diametro $\le h$ \\
\textbf{Group Average} & Media $\frac{\sum \sum d}{|A||B|}$ & Globulare tollerante & Moderata & Bilanciato \\
\textbf{Metodo di Ward} & $\Delta\mathrm{SSE}_{AB}$ & Sferica bilanciata & Molto bassa & Analogo K-Means \\
\bottomrule
\end{tabularx}
\caption{Confronto analitico tra i quattro criteri di collegamento nel clustering gerarchico.}
\label{tab:dm-linkage-criteria-synoptic}
\end{table}

% ====================================================================
% SEZIONE 7.9: CLUSTERING GERARCHICO DIVISIVO BASATO SU MST
% ====================================================================
\section{Clustering Gerarchico Divisivo basato su Minimum Spanning Tree (MST)}
\label{sec:dm-divisive-mst-clustering}

Mentre gli algoritmi divisivi generici risultano computazionalmente proibitivi (considerare tutte le possibili bisezioni di un insieme di $n$ punti comporta $2^{n-1} - 1$ configurazioni), l'approccio basato su teoria dei grafi fornisce una soluzione elegante ed efficiente tramite il \textbf{Minimum Spanning Tree} ($\mathrm{MST}$).

\begin{definition}[Minimum Spanning Tree dei Dati]
Dato il grafo completo $G = (V, E)$ in cui i vertici $V$ rappresentano gli $n$ punti dati e ogni arco $(u, v) \in E$ ha peso $w(u, v) = d(\mathbf{x}_u, \mathbf{x}_v)$, un $\mathrm{MST}$ è un albero ricoprente privo di cicli di peso complessivo minimo:
\begin{equation}
    w(\mathrm{MST}) = \sum_{e \in \mathrm{MST}} w(e) \to \min
\end{equation}
\end{definition}

\subsection{Algoritmo Divisivo basato su Taglio di Archi}

\begin{algorithm}[htbp]
\SetAlgoLined
\KwInput{Dataset $\mathcal{D} = \{\mathbf{x}_1, \dots, \mathbf{x}_n\}$, numero di cluster desiderato $K$.}
\KwOutput{Partizione $\{C_1, \dots, C_K\}$.}
Costruisci il grafo completo pesato sui dati $G = (V, E)$\;
Estrai il Minimum Spanning Tree $\mathcal{T}_{\mathrm{MST}}$ mediante l'algoritmo di Prim o Kruskal\;
Ordina gli $n - 1$ archi dell'albero per peso decrescente: $w(e_1) \ge w(e_2) \ge \dots \ge w(e_{n-1})$\;
\For{$k = 1$ \KwTo $K - 1$}{
    Rimuovi dall'albero l'arco di peso massimo residuo $e_k$\;
    L'albero si scinde in due componenti connesse disgiunte\;
}
Identifica le $K$ componenti connesse risultanti come i $K$ cluster della partizione\;
\caption{Clustering divisivo basato su Minimum Spanning Tree.}
\label{alg:dm-divisive-mst}
\end{algorithm}

\begin{property}[Equivalenza Teorica con il Single Link]
La partizione ottenuta rimuovendo i $K - 1$ archi a peso massimo dal Minimum Spanning Tree è \textbf{identica} alla partizione ottenuta arrestando l'algoritmo agglomerativo con Single Link a $K$ cluster. Ne consegue che la costruzione dell'$\mathrm{MST}$ (realizzabile in $\mathcal{O}(n^2)$ con Prim) fornisce una via computazionalmente diretta per ricavare la gerarchia del Single Link.
\end{property}

% ====================================================================
% SEZIONE 7.10: COMPLESSITÀ, LIMITI E CONFRONTO SINOTTICO
% ====================================================================
\section{Complessità, Limiti e Confronto Sinottico}
\label{sec:dm-hierarchical-complexity-comparison}

\subsection{Analisi della Complessità Asintotica del Clustering Gerarchico}

\begin{itemize}
    \item \textbf{Complessità Spaziale}: L'algoritmo deve memorizzare l'intera matrice di prossimità triangolare $n \times n$:
    \begin{equation}
        \mathcal{S}_{\text{Gerarchico}} = \mathcal{O}(n^2)
    \end{equation}
    Ciò costituisce un limite severo nei contesti di Big Data. Ad esempio, per $n = 100\,000$ punti a precisione standard (8 byte per float), la matrice richiede circa $40\,\text{GB}$ di memoria RAM, rendendo il metodo inapplicabile senza partizionamenti preliminari.
    \item \textbf{Complessità Temporale Naive}: Il ciclo compie $n - 1$ iterazioni. A ogni iterazione, la ricerca lineare del minimo nella matrice di dimensione residua richiede $\mathcal{O}(n^2)$ confronti, portando a:
    \begin{equation}
        \mathcal{T}_{\text{Naive}} = \sum_{t=1}^{n-1} \mathcal{O}((n - t)^2) = \mathcal{O}(n^3)
    \end{equation}
    \item \textbf{Ottimizzazione tramite Code con Priorità (Heap)}: Mantenendo le distanze per ciascun cluster in una coda con priorità (\textit{min-heap}), l'estrazione del minimo e l'aggiornamento richiedono $\mathcal{O}(n \log n)$ per iterazione, riducendo la complessità totale a:
    \begin{equation}
        \mathcal{T}_{\text{Heap}} = \mathcal{O}(n^2 \log n)
    \end{equation}
    (Per criteri specifici come Single Link e Complete Link, algoritmi ottimizzati come SLINK e CLINK operano in $\mathcal{O}(n^2)$ tempo e $\mathcal{O}(n)$ spazio ausiliario).
\end{itemize}

\subsection{Limiti Strutturali del Clustering Gerarchico}

\begin{enumerate}
    \item \textbf{Irreversibilità delle Decisioni (\textit{Greedy Nature})}: Le decisioni di fusione o scissione sono permanenti e prive di backtracking. Se due punti vengono uniti al passo $t$ per via del rumore locale o di un criterio inadeguato, non potranno mai essere separati nei passi successivi.
    \item \textbf{Mancanza di una Funzione Obiettivo Globale}: A eccezione del Metodo di Ward (che minimizza l'incremento locale di varianza), i criteri di linkage non ottimizzano una funzione di costo globale scalare sull'intero dataset.
\end{enumerate}

\subsection{Schema Comparativo Finale: K-Means vs Clustering Gerarchico}

A completamento della trattazione dei due paradigmi, la Tabella~\ref{tab:dm-kmeans-vs-hierarchical} sintetizza le divergenze architetturali, algebriche e operative essenziali che guidano la scelta del modello in contesti ingegneristici e scientifici reali.

\begin{table}[htbp]
\centering
\small
\renewcommand{\arraystretch}{1.3}
\begin{tabularx}{\textwidth}{l Y Y}
\toprule
\textbf{Caratteristica Operativa} & \textbf{K-Means Clustering} & \textbf{Clustering Gerarchico} \\
\midrule
\textbf{Tipo di Struttura} & Partizionale piatta (\textit{hard clustering}) & Tassonomica ad albero nidificato (Dendrogramma) \\
\textbf{Parametro $K$} & Obbligatorio a priori & Derivabile a posteriori tracciando un taglio \\
\textbf{Funzione Obiettivo} & Minimizzazione di $\mathrm{SSE}$ globale & Scelte miopi locali (\textit{greedy}) senza costo globale \\
\textbf{Complessità Temporale} & $\mathcal{O}(n \cdot K \cdot I \cdot d)$ (Lineare con $n$) & $\mathcal{O}(n^2 \log n)$ fino a $\mathcal{O}(n^3)$ (Sovra-quadratica) \\
\textbf{Complessità Spaziale} & $\mathcal{O}((n + K) \cdot d)$ (Lineare con $n$) & $\mathcal{O}(n^2)$ (Quadratica, memoria critica) \\
\textbf{Scalabilità sui Dati} & Estremamente elevata (milioni di record) & Molto bassa ($n \le 20\,000$ in RAM standard) \\
\textbf{Sensibilità Inizializzazione} & Elevatissima (intrappolamento in minimi locali) & Deterministico (nessun seme o partizione iniziale) \\
\textbf{Flessibilità di Forma} & Solo sferica/globulare (iperpiani Voronoi) & Arbitraria con Single Link; sferica con Ward \\
\textbf{Reversibilità Decisioni} & Sì (i punti migrano ad ogni iterazione) & No (fusioni e tagli permanenti senza riassegnazione) \\
\bottomrule
\end{tabularx}
\caption{Confronto sinottico globale tra K-Means e Clustering Gerarchico.}
\label{tab:dm-kmeans-vs-hierarchical}
\end{table}
